%======================================================================
%  The processor: Intel Xeon CPU Max 9480
%======================================================================

\section{The Xeon CPU Max 9480 processor}
\label{sec:xeon-max-processor}

The Xeon CPU Max series is the Sapphire Rapids generation of Xeon
Scalable processors with high-bandwidth memory added to the package.
The Max 9480 installed on CRESCO8 is the series' largest part: 56 cores,
112.5\,MB of last-level cache, 64\,GB of on-package HBM2e, eight DDR5-4800
channels and four UPI links at 16\,GT/s, with a 1.9\,GHz base frequency,
a 3.5\,GHz maximum turbo frequency and a 350\,W thermal design
power~\cite{intel9480,intelmaxoverview}.  This section describes the
parts of the processor that determine how an inference engine's work
and data map onto it; the memory systems follow in
\secref{sec:xeon-max-platform}.

\subsection{Tiles, mesh and locality}
\label{sec:xeon-max-tiles}

A Max 9480 package is not one monolithic die.  It contains four compute
tiles joined by EMIB silicon bridges, which extend the on-die mesh across
tile boundaries so that the four tiles behave as one cache-coherent
processor~\cite{nassif2022sapphire}.  On the Max 9480 each tile enables 14
cores, holds a share of the last-level cache, and carries one DDR5 memory
controller with two channels and one 16\,GB HBM2e
stack~\cite{intelmaxoverview,lp1738}.  The memory controllers reported by
the Linux EDAC driver on a CRESCO8 HBM node match this arrangement: four
DDR5 controllers with two channels each per socket, and HBM controllers
whose pseudo-channels add up to 64\,GB per socket.

\figref{fig:node-locality} draws the resulting node.  The tile is the
physical origin of the locality structure discussed later.  With
sub-NUMA clustering disabled, addresses are interleaved over all four
tiles and the socket appears as one locality domain.  With SNC4, each
tile together with its two DDR5 channels and its HBM stack becomes a
separate NUMA domain (\secref{sec:hbm-clustering}).  Two sockets are
connected by UPI, so the dual-socket node is one shared-memory system
with two levels of distance: tile to tile within a package, and package
to package across UPI.

\begin{figure}[htbp]
\centering
\begingroup
\definecolor{socketink}{HTML}{263445}
\definecolor{socketblue}{HTML}{365D85}
\colorlet{socketaccent}{sapred}
\definecolor{socketborder}{HTML}{CBD5DF}
\resizebox{\textwidth}{!}{%
\begin{tikzpicture}[
  font=\small\sffamily, text=socketink,
  card/.style={line width=0.65pt, rounded corners=4pt, align=center,
               inner sep=3pt},
  tile/.style={card, draw=socketborder, fill=socketink!4,
               minimum width=2.6cm, minimum height=1.55cm},
  hbm/.style={card, draw=socketaccent, fill=socketaccent!9,
              minimum width=1.22cm, minimum height=0.82cm,
              font=\footnotesize\sffamily},
  ddr/.style={card, draw=socketblue, fill=socketblue!9,
              minimum width=1.22cm, minimum height=0.82cm,
              font=\footnotesize\sffamily},
  bridge/.style={line width=2.2pt, draw=socketink!35},
  wire/.style={line width=0.8pt},
  note/.style={font=\footnotesize\sffamily, text=socketink!80,
               align=center},
  link/.style={{Stealth[length=5pt,width=4pt]}-{Stealth[length=5pt,width=4pt]},
               line width=0.8pt, draw=socketink!75}
]
  \foreach \s/\x in {0/-4.0, 1/4.0} {
    % Four tiles in a 2x2 arrangement; the grid conveys the tile
    % structure, not the exact floorplan.
    \foreach \t/\dx/\dy in {a/-1.55/0.95, b/1.55/0.95, c/-1.55/-0.95, d/1.55/-0.95} {
      \node[tile] (t\s\t) at ({\x+\dx},\dy)
        {\textbf{Tile}\\14 cores\\{\footnotesize L2 + LLC slices}};
    }
    \draw[bridge] (t\s a.east) -- (t\s b.west);
    \draw[bridge] (t\s c.east) -- (t\s d.west);
    \draw[bridge] (t\s a.south) -- (t\s c.north);
    \draw[bridge] (t\s b.south) -- (t\s d.north);
    \node[draw=socketborder, line width=0.65pt, rounded corners=6pt,
          inner sep=6pt, fit=(t\s a)(t\s d)] (pkg\s) {};
    % Each tile owns one HBM2e stack and two DDR5 channels.
    \foreach \t/\dir/\anc in {a/above/north, b/above/north, c/below/south, d/below/south} {
      \node[hbm, \dir=0.55cm of t\s\t.\anc, xshift=-0.66cm] (h\s\t)
        {\textbf{HBM2e}\\16\,GB};
      \node[ddr, \dir=0.55cm of t\s\t.\anc, xshift=0.66cm] (d\s\t)
        {\textbf{DDR5}\\2 ch};
      \draw[wire, draw=socketaccent] (h\s\t) -- (h\s\t |- t\s\t.\anc);
      \draw[wire, draw=socketblue] (d\s\t) -- (d\s\t |- t\s\t.\anc);
    }
    \node[align=center, below=0.15cm of h\s c.south -| pkg\s.south]
      {\textbf{Socket \s}\\
       {\footnotesize 56 cores $\cdot$ 4 $\times$ 16\,GB HBM2e $\cdot$
        8 DDR5 channels}};
  }
  \draw[link] (pkg0.east) -- node[above, note]{UPI}
    node[below, note]{4 links\\16\,GT/s} (pkg1.west);

  % One SNC4 domain: a tile with its own channels and HBM stack.
  \node[draw=socketaccent!80, densely dashed, line width=0.7pt,
        rounded corners=5pt, inner sep=3pt,
        fit=(t1b)(h1b)(d1b)] (snc) {};
  \node[note, right=0.15cm of snc.east, anchor=west, text=socketaccent]
    {one SNC4\\domain};
\end{tikzpicture}%
}
\endgroup
\caption{Tile-level locality map of a CRESCO8 Xeon Max node.  Each
package contains four compute tiles bridged by EMIB (inter-tile bars); every
tile owns two DDR5 channels and one HBM2e stack.  The dashed outline
marks what becomes one NUMA domain under SNC4.  The drawing shows
ownership and locality, not the physical floorplan~\cite{nassif2022sapphire,intelmaxoverview,lp1738}.}
\label{fig:node-locality}
\end{figure}

\subsection{Cores and caches}
\label{sec:xeon-max-cores}

The cores are Golden Cove server cores.  On the allocated node
\texttt{lscpu} reports a 48\,KB, 12-way L1 data cache and a 2\,MB,
16-way private L2 per core, and a 112.5\,MB, 15-way last-level cache
shared by the 56 cores of a socket.  The last-level cache is
non-inclusive: it is not guaranteed to hold a copy of the lines in the L2
caches~\cite{intelmaxoverview}.

Two properties of this hierarchy matter for inference.  First, the
whole last-level cache is smaller than the weights of any model studied
here, so the weights of every decode step stream from memory rather than
from cache.  Second, the number of misses a core can keep outstanding is
bounded.  Each Golden Cove core has 16 L1 data-cache fill buffers and can
track 48 outstanding misses in its memory subsystem~\cite{anandtech2021goldencove}.  These
limits set how much memory-level parallelism one core can generate.
Together with Little's law (\secref{sec:littles-law}) they determine how
many cores are needed before a memory tier's bandwidth becomes reachable
(\secref{sec:platform-concurrency}).

\subsection{Vector and matrix instructions}
\label{sec:xeon-max-isa}

Each core has two 512-bit fused multiply-add units for AVX-512.  On top
of the floating-point instructions, the core supports AVX-512 VNNI for
8-bit integer dot products, AVX-512 BF16 and AVX-512 FP16.  The largest
addition is Advanced Matrix Extensions (AMX).  AMX adds eight tile
registers of 1\,KB each, organised as up to 16 rows of 64 bytes, and a tile
matrix-multiply unit that executes BF16 and INT8 multiply-accumulate on
whole tiles.  Intel gives its peak as 1{,}024 BF16 and 2{,}048 INT8
operations per cycle per core, against 64 and 256 for the AVX-512 BF16 and
VNNI paths~\cite{intelarchday2021}.  The node's \texttt{/proc/cpuinfo}
lists the \texttt{amx\_tile}, \texttt{amx\_bf16} and \texttt{amx\_int8}
flags together with \texttt{avx512\_bf16}, \texttt{avx512\_fp16} and
\texttt{avx512\_vnni}.

AMX is the reason that matrix-matrix work such as prompt prefill is
practical on this CPU: oneDNN, PyTorch and the CPU backends of
\lcpp{} and vLLM dispatch BF16 and INT8 GEMMs to it.  It does little for
low-batch decode, whose matrix-vector products are limited by the rate at
which weights arrive rather than by arithmetic (\secref{sec:platform-llm}).
AMX also has an operating-system requirement.  Its 8\,KB of tile state is
enabled per process: a program must request permission with
\texttt{arch\_prctl(ARCH\_REQ\_XCOMP\_PERM)} before its first AMX
instruction~\cite{linuxxstate}.  The node kernel advertises the AMX
flags, so the mechanism is available; the inference runtimes make the
request themselves.

\subsection{Clocks, threads and nominal peaks}
\label{sec:xeon-max-peaks}

The 1.9\,GHz base and 3.5\,GHz turbo frequencies are not the frequencies
at which dense vector code runs.  Heavy AVX-512 and AMX execution lowers
the sustainable core clock, and an all-core load is further limited by
the 350\,W package budget.  Peak rates computed at the base frequency are
therefore an upper bound for sustained vector work, not a prediction.
On CRESCO8, simultaneous multithreading is not active: each node exposes
112 logical CPUs, one per physical core, and Slurm schedules them with one
thread per core.  The operating system manages frequency through the
\texttt{acpi-cpufreq} driver with the \texttt{schedutil} governor,
because the kernel is booted with \texttt{intel\_pstate=disable}
(\secref{sec:cresco8-observed}).

\tabref{tab:xeon-max-peaks} collects the nominal per-socket peaks used in
the rest of the thesis.  The compute rows multiply the per-core rate by 56
cores and the base frequency.  The DDR5 row is the channel arithmetic,
eight channels of 4{,}800\,MT/s with eight bytes per transfer.  Intel does
not publish a channel-level HBM figure and quotes up to about 1\,TB/s per
socket~\cite{intelmaxoverview}.  These are specification values.  The
sustainable ceilings, measured with dedicated bandwidth and compute
kernels and combined into measured Rooflines, are the subject of the next
chapter.

\begin{table}[htbp]
\centering
\small
\tabsetup
\begin{tabularx}{\textwidth}{P{3.4cm}P{2.7cm}P{2.4cm}Y}
\headrow
\headcell{Resource} & \headcell{Per-core rate} & \headcell{Per socket} &
\headcell{Basis} \\[3pt]
\tabheadrule
\tabstripeRow
AVX-512 FP64 & 32 FLOP/cycle & 3.4\,TFLOP/s & Two FMA units, 8 lanes \\[3pt]

\tabrowsep
\tabplainRow
AVX-512 FP32 & 64 FLOP/cycle & 6.8\,TFLOP/s & Two FMA units, 16 lanes \\[3pt]

\tabrowsep
\tabstripeRow
AMX BF16 & 1{,}024 FLOP/cycle & 108.9\,TFLOP/s & Intel peak
  figure~\cite{intelarchday2021} \\[3pt]

\tabrowsep
\tabplainRow
AMX INT8 & 2{,}048 op/cycle & 217.9\,TOP/s & Intel peak
  figure~\cite{intelarchday2021} \\[3pt]

\tabrowsep
\tabstripeRow
DDR5-4800, 8 channels & -- & 307.2\,GB/s & $8 \times 4{,}800\,\text{MT/s}
  \times 8$\,B \\[3pt]

\tabrowsep
\tabplainRow
HBM2e, 4 stacks & -- & ${\approx}1$\,TB/s & Intel's quoted
  figure~\cite{intelmaxoverview} \\[3pt]

\tabrowsep
\tabstripeRow
UPI 2.0 & -- & 4 links, 16\,GT/s & Socket-to-socket
  links~\cite{intel9480} \\[3pt]
\tabbottomrule
\end{tabularx}
\caption{Nominal per-socket peaks of the Max 9480.  Compute rates assume
all 56 cores at the 1.9\,GHz base frequency and are upper bounds; the
memory rows are interface peaks, not sustainable bandwidths.}
\label{tab:xeon-max-peaks}
\end{table}
